\songtitle{Señora Rosario's House}

\tag*{2026}\tag{folk}\tag{ricochets}\tag{rosario-cubides}\tag{spanish-lyrics}
\begin{notes}
From Carlos Thompson's novel \emph{Ricochets}. Character profile of Rosario Cubides by Claude. Lyrics by Suno.
\end{notes}
\begin{style}
folk, pasillo texture, latin jazz instrumental bridges, spoken word fragments, mid tempo 98 bpm, male baritone lead vocals, middle-aged female spoken voice, central Colombian accent, nylon-string guitar arpeggios, acoustic bass, brushed snare, conga rim clicks, trumpet stabs, Hammond organ pads, tape saturation, plate reverb, 2000s stereo mix, street ambience, bittersweet nostalgia, domestic realism
\end{style}
\begin{lyrics}
\songpart{Intro}
\annotation{instrumental — guitar, a street sound fading in: vendors, a bus brake,}
\annotation{male baritone vocals}\\
Second floor, old colonial stone,\\
a house on a street that was never his own.

\annotation{Spoken — female, central Colombian, fast-paced, shouted}
``Ya no aguanto al gringo ese... no hace sino llegar borracho a cualquier hora.''

\songpart{Verse 1 — male baritone vocals}
He opens his eyes on the hallway floor,\\
a boy's small hands pull him up off the ground.\\
Curtains stay shut, let the daylight wait,\\
a phone with no time, a bottle he throws down.\\
He comes down the stairs when the smell has cleared,\\
sits at the table where the flan dish waits.\\
She looks at the boy, and she looks at him,\\
and nobody needs to translate that face.

\annotation{Spoken — female, central Colombian, fast-paced, whispered}
``¿Se lo anotaste a la cuenta?''

\songpart{Chorus — male baritone vocals}
Rosario keeps the ledger and the door,\\
Camilo runs the errands, Gina holds the floor,\\
but it's her house, her rules, her roof, her name —\\
the gringo drinks, the gringo comes home late,\\
and every single one of us knows whose house this stays.

\songpart{Verse 2 — male baritone vocals}
Windows Server humming on a borrowed machine,\\
Wimbledon, the Tour, a race gone wrong in France —\\
Érika, Lida, Camilo at the screen,\\
till the beans and the bacon cut the debate short.

\annotation{Spoken — female, central Colombian, fast-paced, shouted}
``Muchachos, el almuerzo está listo... Camilo, ayúdeme a servir.''

\songpart{Chorus — repeat, slight shift}
Rosario keeps the schedule and the door,\\
Camilo does the serving, Gina keeps the score,\\
but it's her house, her table, her name —\\
the fríjoles are calling, the guesswork can wait,\\
and every single one of us comes running just the same.

\songpart{Bridge — male baritone vocals}
One more suitcase zipped by the door,\\
one more girl who won't live here anymore.\\
Hands find hands across the kitchen light,\\
nothing left to say gets said just right.\\
A boy taps his chest, she mirrors the same,\\
a whole year of living folds into a name.

\annotation{Spoken — female, central Colombian, almost crying}
``Vete antes de que me hagas llorar.''

\songpart{Final Chorus — male baritone vocals}
Rosario keeps the house and the door,\\
some come, some go, but she's here like before,\\
it was never really theirs, it was always her name —\\
the ones who left it changed, the ones who stayed the same,\\
and every single one of them still calls this place home.

\songpart{Outro — Spoken — female, central Colombian, dismisive}
``Deje así.''
\end{lyrics}

